\section{The algebraic consequence of the entropy trace identity}
\label{alg:section}

We work over $\C$, by complexifying the real formal symmetry algebra.
Real structures will be retained throughout and recovered at the end.
The notation and sign convention are those of Section~\ref{sym:section}.
In particular, set
\[
 E=\mathfrak h/\mathfrak k,
 \qquad
 \omega(p+\mathfrak k,q+\mathfrak k)=\xi([p,q]).
\]
An overall change of sign in this identification has no effect below.
We use the following conclusions of Propositions~\ref{sym:dilation}
and~\ref{sym:effective} and Lemma~\ref{sym:trace}:
\begin{enumerate}[label=(\roman*)]
\item $\mathfrak h$ is a linearly compact Lie algebra with continuous
bracket, and $\mathfrak k$ is an open subalgebra containing no nonzero
ideal of $\mathfrak h$.
\item The continuous functional $\xi$ has coadjoint stabilizer
$\mathfrak k$, and the displayed form on the finite-dimensional space
$E$ is symplectic.
\item $H\in\mathfrak k$ acts by a real, degreewise semisimple formal
grading. Its adjoint powers preserve the ambient finite-jet kernels.
Closed invariant spaces admit continuous pure-weight projections, and
zero-weight vectors in finite-dimensional continuous invariant quotients
admit zero-weight lifts.
\item $\mathfrak h_0\subset\mathfrak k$, and $E$ has no zero
$H$-weight. For a fixed nonzero real constant $c$,
\begin{equation}
 c\,\xi(p)=\tr\bigl(p|E_+\bigr),
 \qquad p\in\mathfrak h_0.
 \label{alg:entropy-trace}
\end{equation}
Here the action on $E$ is the bracket isotropy action, and $+$ denotes
positive $\ad H$-weights.
\end{enumerate}
The grading on $E$ has signs opposite to the eigenvalues of the projected
dilation $D$. Thus its positive part corresponds to the actual stable
space. We use algebraic signs in the calculations and return to the
actual stable and unstable notation when constructing fields.

We seek invariant formal directions of each dynamical sign. The
construction has two algebraic steps. First we find an intrinsic
coisotropic subspace of $E$; its symplectic orthogonal will provide the
characteristic directions in the infinite-dimensional case. Then the
trace identity forces the sign subspaces to belong to finite intrinsic
lists. That finiteness will let us transport the choices smoothly along
actual axes, while their ambient extensions remain formal.

\subsection{The intrinsic coisotropic reduction}

For $p\in\mathfrak h$, consider the family of operators
$\{(\ad p)^n:n\ge0\}$. Equicontinuity in the formal linear topology
means that, for each prescribed output jet order, there is an input
vanishing order that makes every operator in this family have zero
output jet. The input order may depend on the output order and on $p$,
but not on $n$. Define the set
\begin{equation}
 \mathfrak l=\{p\in\mathfrak h:
       \{(\ad p)^n:n\ge0\}\text{ is equicontinuous}\}.
 \label{alg:equicontinuity}
\end{equation}
In finite dimension this is all of $\mathfrak h$, because the topology
is discrete. In general, even its subalgebra property needs proof.

\begin{proposition}[Intrinsic coisotropic reduction]
\label{alg:coisotropic}
The set $\mathfrak l$ is an open intrinsic subalgebra containing
$\mathfrak k$ and $H$, and has arbitrarily small open ideals for its
induced topology. The subspace $L=\mathfrak l/\mathfrak k$ of $E$ is
coisotropic, meaning $L^\omega\subset L$. If $\mathfrak h$ is infinite
dimensional, then $\mathfrak l\ne\mathfrak h$.
\end{proposition}

We will identify $\mathfrak l$ as an origin stabilizer in the classical
realizations of the minimal ideals of $\mathfrak h$. The trace identity
will then show that $\mathfrak k$ fixes those origins and prove
coisotropy. We begin by finding the ideals on which to make this test.

\begin{lemma}
\label{alg:minimal-ideals}
The algebra $\mathfrak h$ has no nonzero closed abelian ideal. Every
nonzero closed ideal contains a minimal nonzero closed ideal. There
are finitely many such minimal ideals, each is topologically perfect,
and their common centralizer in $\mathfrak h$ is zero.
\end{lemma}

\begin{proof}
Let $A$ be a closed abelian ideal. For $p\in A\cap\mathfrak h_0$ we
have $p\in\mathfrak k$ and $(\ad p)^2=0$ on $\mathfrak h$.
Consequently its induced trace on $E_+$ vanishes, and
\eqref{alg:entropy-trace} gives $\xi(p)=0$. The functional $\xi$ kills
every nonzero weight because $H$ stabilizes it. Weight projection
therefore gives $\xi(A)=0$. Since $[\mathfrak h,A]\subset A$, every
element of $A$ stabilizes $\xi$; hence $A\subset\mathfrak k$, which
forces $A=0$.

Consider a decreasing chain of nonzero closed ideals contained in a
given nonzero closed ideal. Their images in $E$ are nonzero: an ideal
with zero image would be contained in $\mathfrak k$. These images are
nested finite-dimensional spaces. Choose a member whose image has
minimal dimension, and a nonzero vector in that image. On the cofinal
subchain below this member, the fibers over this vector form a nested
family of nonempty closed affine subspaces of a linearly compact
space. Their intersection is nonempty. The intersection of the chain
is consequently a nonzero closed ideal. The usual minimal-element
argument now gives a minimal nonzero closed ideal.

For a minimal ideal $I$, the closure of $[I,I]$ is a closed ideal of
$\mathfrak h$. It is nonzero by the first paragraph, so it equals $I$.
Also $\xi(I)\ne0$: otherwise the ideal property would put $I$ in
$\mathfrak k$. By continuity and topological perfection, there exist
$a,b\in I$ such that $\xi([a,b])\ne0$. Their values in $E$ span a
symplectic two-plane.

Distinct minimal ideals commute. Indeed their bracket is contained
in their intersection, and a nonzero intersection would make them
equal. The two-planes just constructed for distinct ideals are thus
mutually symplectically orthogonal. There are at most $\dim(E)/2$ of
them. Finally, the common centralizer is a closed ideal. If nonzero,
it contains a minimal ideal; this ideal would commute with itself,
contradicting the first paragraph.
\end{proof}

We recall precisely the classical structure results used here. Over
$\C$, the Cartan--Guillemin theorem identifies a nonabelian minimal
closed ideal of a transitive linearly compact Lie algebra with
\begin{equation}
 I=S\widehat\otimes\mathcal O_e,
 \qquad
 \mathcal O_e=\C[[z_1,\ldots,z_e]],
 \label{alg:current-ideal}
\end{equation}
where $S$ is a nonabelian simple linearly compact Lie algebra and $e$
is finite. This is a topological Lie algebra identification, with
pointwise current bracket. In the coinduced formulation, a maximal
proper closed ideal of $I$ has an open normalizer $N$ in $\mathfrak h$;
the simple quotient is $S$, and
$I\simeq\Hom_{U(N)}(U(\mathfrak h),S)$. The finite complement to $N$
and the PBW coalgebra identification give~\eqref{alg:current-ideal}.
See \citet[Proposition~2.6, Theorem~2.7 and Corollary~2.8]{FattoriKac2002}
and \citet[Proposition~6.11]{BakalovDAndreaKac2001}.

The infinite-dimensional possibilities for $S$ are the formal Cartan
algebras: all vector fields $W$, divergence-free vector fields,
Hamiltonian vector fields, and contact vector fields. The
finite-dimensional simple algebras are also allowed. For the infinite
simple factors, continuous derivations are represented faithfully by
their normalizing formal vector fields. For the divergence-free and Hamiltonian cases these
normalizers allow a constant conformal multiplier; the normalizers of
$W$ and of the contact algebra are those algebras themselves. These
are the classification and derivation results in
\citet[Theorem~6.8 and Proposition~6.12(i)]{BakalovDAndreaKac2001}.
For a finite-dimensional simple factor, derivations are inner and
are treated as finite matrices.
The topological Lie algebra statements use formal linear topologies
with discrete scalar field, not analytic convergence topologies.

For later order tests, give ordinary coordinates weight one in the
$W$, divergence-free and Hamiltonian realizations. Vector-field grade
is coefficient degree minus one; for a Hamiltonian it is Hamiltonian
degree minus two, modulo constants. In contact coordinates the
horizontal coordinates have weight one and the Reeb coordinate has
weight two; contact-Hamiltonian grade is weighted degree minus two.
Write $S_{\ge j}$ for series with all grades at least $j$, and
$S_{(0)}$ for the subalgebra fixing the origin. In these realizations,
$S_{(0)}=S_{\ge0}$. The lowest grades are $-1$, except in the contact
case, where they are $-2$ and $-1$.

\begin{lemma}
\label{alg:derivation-realization}
For an ideal~\eqref{alg:current-ideal}, every continuous derivation has
a unique decomposition into a formal parameter vector field and an
$\mathcal O_e$-family of continuous derivations of $S$:
\begin{equation}
 \Der(I)=W_e\otimes1+
 \mathcal O_e\widehat\otimes\Der(S).
 \label{alg:current-derivations}
\end{equation}
The action of $\mathfrak h$ on its finitely many minimal ideals gives
a faithful topological embedding with closed image in the product of
these derivation algebras.
\end{lemma}

\begin{proof}
The continuous centroid of a simple factor $S$ is scalar. One may
use the linearly compact Schur lemma, as in
\citet[Lemma~2.2]{FattoriKac2002}. In the listed realizations it also
follows by commuting a centroid map with the grading in the $W$ and
contact cases, or with the degree-zero $\mathfrak{sl}$ or
$\mathfrak{sp}$ action in the other cases. A grading eigenspace in
the former cases, and a finite-multiplicity isotypic component in
the latter, is a nonzero finite-dimensional centroid-invariant
space. Such components exist because the homogeneous
divergence-free polynomial fields and homogeneous Hamiltonians
form irreducible modules with distinct highest weights as their
degrees vary. An eigenvalue on this finite-dimensional space gives
a nonzero kernel for the centroid map minus that scalar. This
kernel is an ideal, so simplicity makes the map scalar.

A continuous centroid map of $I$, restricted to constant $S$-series
and tested coefficientwise against brackets with constants, is
multiplication by a scalar series $a(z)$. To see that this determines
the map everywhere, subtract this multiplication map. For each
monomial $z^\nu$ its value on $z^\nu S$ is zero: write elements of
$S$ as limits of finite sums of brackets and use
$z^\nu[a,b]=[z^\nu a,b]$, placing the centroid map on the constant
entry. Polynomial series are dense, so continuity finishes the
argument. Thus the centroid is exactly $\mathcal O_e$.

A derivation acts on this centroid by commutator and hence induces a
continuous derivation of $\mathcal O_e$, namely a formal parameter
vector field. For continuity, evaluate the commutator on a fixed
nonzero constant $s\in S$ and apply a continuous linear functional
on $S$ taking value one on $s$. Continuity on $I$ then gives
continuity of the induced scalar-series operator. After subtracting
the parameter field, the remaining derivation is
$\mathcal O_e$-linear. Restriction to constant $S$-series, coefficient
by coefficient, gives derivations of $S$. Density again gives the
decomposition and its uniqueness. This proves
\eqref{alg:current-derivations}; it is also the continuous-derivation
formula of \citet[Proposition~6.12(ii)]{BakalovDAndreaKac2001}.

The kernel of the action on all minimal ideals is their common
centralizer, which is zero by Lemma~\ref{alg:minimal-ideals}.
The displayed derivation spaces are linearly compact in their formal
coefficient topology: factor normalizers are closed spaces of formal
vector fields, finite simple factors give finite matrix spaces, and
completed coefficient products preserve linear compactness.
Continuity in the asserted jet topologies follows from continuity of
bracketing. More explicitly, the normalizer realization embeds
continuously into the product of its bracket evaluations on factor
elements. It is injective and has continuous inverse on its image:
a continuous injective map from a linearly compact space to a
Hausdorff linearly topologized space has closed image and induces its
original topology. Apply this fact first to the normalizer and then
to the faithful map from $\mathfrak h$.
\end{proof}

\paragraph{Detecting translations.}

In the faithful realization just obtained, fields fixing the origins
preserve high-order jet kernels. We next show that a translation has
the opposite behavior, so that the origin stabilizers can be identified
with $\mathfrak l$.

\begin{lemma}
\label{alg:translation-test}
Let $A$ be a normalizing formal vector field of an infinite simple
Cartan algebra, and suppose that it does not fix the origin. Its
lowest nonzero grade is $-b<0$. The grade $-b$ part of $\ad A$ maps
each homogeneous component of grade $j+b$ onto the component of
grade $j$. Consequently the family $\{(\ad A)^n:n\ge0\}$ is not
equicontinuous for the formal linear topology.
\end{lemma}

\begin{proof}
For $W$ this is surjectivity of differentiation in a nonzero constant
direction on homogeneous vector-field coefficients. For
divergence-free fields, integrate the coefficients in that direction.
The divergence of the resulting primitive is independent of this
coordinate. Integrating that error in another coordinate and
subtracting the resulting component corrects the divergence without
changing the first directional derivative. The correction is
homogeneous. In the Hamiltonian case, integrate homogeneous
Hamiltonians, modulo constants.

For contact fields, a grade $-2$ Hamiltonian acts by Reeb
differentiation. A nonzero grade $-1$ Hamiltonian acts by a constant
horizontal derivative plus a linear multiple of Reeb differentiation.
Its Reeb derivative is zero, so the contact bracket on Hamiltonians
has no additional multiplier term. A homogeneous quadratic change
of the Reeb coordinate straightens this directional operator.
Integration in the resulting coordinate preserves the indicated
weighted degrees and proves surjectivity.

Fix a grade $j$ with a nonzero homogeneous element. Repeatedly choose
homogeneous primitives to obtain elements of grade $j+nb$ whose
$n$th image under the leading operation is that fixed element.
In $(\ad A)^n$, only the product of the lowest-grade operations can
reach grade $j$ from grade $j+nb$; all other products have larger
grade. The inputs tend to zero in the formal topology, while these
outputs do not. This contradicts equicontinuity.
\end{proof}

For each infinite-dimensional minimal ideal $I_i$, let
$\mathfrak l_i$ be the inverse image in $\mathfrak h$ of the stabilizer
of the parameter origin and, when its simple factor is Cartan, the
factor origin at zero parameters. Finite-dimensional minimal ideals
impose no condition. These are subalgebras defined by finitely many
coefficient conditions. Their intersection gives the following concrete
description of the intrinsic set just defined.

\begin{lemma}[Origin stabilizers]
\label{alg:intrinsic-l}
The set $\mathfrak l$ is the intersection of the subalgebras
$\mathfrak l_i$. In particular it is open and intrinsic, contains $H$,
and has arbitrarily small open ideals for its induced topology.
\end{lemma}

\begin{proof}
There are finitely many origin conditions, each involving finitely
many coefficients, so their intersection is open. In the faithful
realization of Lemma~\ref{alg:derivation-realization}, an element of
this intersection acts by vector fields fixing the joint origin. Its
brackets preserve arbitrarily high total-order jet kernels. In joint
parameter and factor coordinates, this is ordinary coefficient order:
differentiation loses one order and multiplication by an
origin-vanishing coefficient restores it. In the
case of a finite simple factor with parameters, the same assertion
holds for a parameter field fixing zero and a matrix-valued series.
On each finite-dimensional ideal the discrete topology imposes no
additional restriction. The faithful topological embedding proves
equicontinuity on $\mathfrak h$.

Conversely, suppose the parameter field of $p$ has a nonzero value
at zero. Choose $z_1$ with $p(z_1)(0)\ne0$ and a fixed nonzero
$s\in S$. Applied $n$ times to $z_1^n s$, only repeated use of the
constant parameter derivative can produce parameter order zero in
$n$ steps. Its value there is
$n!\,p(z_1)(0)^n s\ne0$. Thus the adjoint powers fail
equicontinuity already on that ideal. If the parameter origin is
fixed but the factor origin is not, reduce modulo the parameter
ideal and apply Lemma~\ref{alg:translation-test}. Equicontinuity
would restrict to the closed invariant ideal $I$ and descend through
the open quotient map $I\to I/(z_1,\ldots,z_e)I\simeq S$, since
the parameter ideal is then invariant. The same
failure holds on $\mathfrak h$ with its induced topology. This
identifies the intersection with the set in
\eqref{alg:equicontinuity}; its intrinsic character follows from that
definition.

The actual formal grading of $H$ preserves finite-jet kernels, so
$H\in\mathfrak l$. Finally, inside $\mathfrak l$ impose vanishing
of all sufficiently low jets in its infinite-ideal realizations and
zero action on its finite-dimensional ideals. These are open ideals
of $\mathfrak l$: bracketing with an origin-preserving field
preserves the prescribed order. They form a neighborhood basis of
zero by the topological embedding.
\end{proof}

\begin{proof}[Proof of Proposition~\ref{alg:coisotropic}]
Lemma~\ref{alg:intrinsic-l} proves the topological assertions.
Only the infinite case needs a further argument. Fix an infinite minimal
ideal $I=S\widehat\otimes\mathcal O_e$, and let
$\mathfrak m=(z_1,\ldots,z_e)$. The space $\mathfrak m I$ is
$H$-invariant because $H$ fixes the parameter origin. If
$p\in\mathfrak m I$ has weight zero, it belongs to
$\mathfrak k$. Its first bracket with $\mathfrak h$ lies in $I$;
each further bracket with $p$ raises parameter order. Since
$\mathfrak k$ is open, this induced action is nilpotent on $E$.
Equation~\eqref{alg:entropy-trace} kills $\xi(p)$. Weight projection
and the vanishing of $\xi$ on nonzero weights give
\begin{equation}
 \xi(\mathfrak m I)=0.
 \label{alg:parameter-annihilation}
\end{equation}
The resulting functional on constant $S$-series is nonzero.

For $q\in\mathfrak k$, test $\xi([q,\mathfrak m I])=0$ on
$z_a s$, choosing $s$ with $\xi(s)\ne0$. At parameter order zero
this bracket is $q(z_a)(0)s$. Thus the parameter field of $q$ fixes
zero. If $S$ is Cartan, let $j$ be the highest standard grade on
which the functional on $S$ is nonzero; it exists by continuity.
If the factor field of $q$ does not fix the origin, let its leading
grade be $-b$. Test on grade $j+b$. By
Lemma~\ref{alg:translation-test} the leading term can have any
prescribed grade-$j$ value, whereas all higher terms have grades
annihilated by $\xi$. This contradicts $q\in\mathfrak k$.
Therefore $\mathfrak k\subset\mathfrak l_i$ for every $i$.

In the contact case $j=-2$ is impossible. Indeed the functional would
then vanish on $\mathfrak m I+S_{(0)}$ and descend nontrivially to
$I/(I\cap\mathfrak l)$. This quotient is exactly
$S/S_{(0)}$: the ideal acts trivially on all the other minimal
ideals, and its factor action is inner. It is an $H$-invariant
subspace of $\mathfrak h/\mathfrak l$. The descended functional is
$H$-invariant because $\xi([H,I])=0$, forcing a zero weight in this
quotient of $E$. But $E$ has no zero weight. Thus $j\ge-1$. In contact
coordinates this also says that a functional supported only on the
Reeb value cannot occur.

For this ideal define the closed linear subspace
\[
 J_i=\mathfrak m I+S_{\ge j+1},
\]
omitting the second term when $S$ is finite dimensional. It lies in
$\mathfrak l_i$, since its possible Cartan term has nonnegative
grade, and in all the other origin stabilizers by commutation of
distinct ideals. Hence $J_i\subset\mathfrak l$.

We claim that the annihilator in $E$ of the image of $J_i$ is
$\mathfrak l_i/\mathfrak k$. If $q\in\mathfrak l_i$, order
preservation at the origins and~\eqref{alg:parameter-annihilation}
give $\xi([q,J_i])=0$. Conversely, this vanishing tested on
$z_a s$ detects every parameter translation. Once those vanish,
the grade-$j+b$ test from the preceding paragraph detects any factor
translation. This proves the claim. Taking symplectic annihilators
and using finite dimensionality of $E$ gives
\[
 L^\omega
 =\left(\bigcap_i\mathfrak l_i/\mathfrak k\right)^\omega
 =\sum_i (\mathfrak l_i/\mathfrak k)^\omega
 =\sum_i\operatorname{image}(J_i)\subset L.
\]

A Cartan factor makes $\mathfrak l$ proper because its inner factor
translations survive modulo $\mathfrak l$. If there are only finite
simple factors but a nontrivial parameter ring, equality
$\mathfrak l=\mathfrak h$ would make $\mathfrak m I$ a nonzero
closed ideal of $\mathfrak h$ killed by $\xi$. It would be
contained in $\mathfrak k$, a contradiction. Without a Cartan
factor or parameters, the faithful derivation realization is finite
dimensional. Thus $\mathfrak l\ne\mathfrak h$ whenever $\mathfrak h$
is infinite dimensional.

\end{proof}

The characteristic space of $L$ is its symplectic orthogonal. To use it
in the trace argument, we need a representation on that space extending
the isotropy action. Set
\begin{equation}
 B=\mathfrak h/\mathfrak l,
 \qquad P=L^\omega,
 \qquad
 \mathfrak u=\operatorname{stab}_{\mathfrak l}(\xi|_{\mathfrak l}).
 \label{alg:base-characteristic}
\end{equation}

\begin{lemma}[The characteristic representation]
\label{alg:characteristic-representation}
We have $\mathfrak u/\mathfrak k=P$, and symplectic pairing identifies
$P$ with $B^*$ as a $\mathfrak k$-module. Transferring the dual bracket
representation on $B^*$ gives a representation of all of $\mathfrak l$
on $P$ extending this $\mathfrak k$ action.
\end{lemma}

\begin{proof}
The kernel of the restricted alternating
form on $L$ is $\mathfrak u/\mathfrak k$, so it is $P$.
The map
\[
 P\longrightarrow B^*,\qquad
 v\longmapsto\bigl(q+\mathfrak l\longmapsto
 \omega(v,q+\mathfrak k)\bigr)
\]
is well-defined and an isomorphism by nondegeneracy. The
$\mathfrak k$ action is symplectic, since its elements stabilize
$\xi$, so this map intertwines its action with the dual of the
bracket action on $B$. The latter is an actual representation of
all of $\mathfrak l$, because $\mathfrak l$ is a subalgebra.
Transferring that representation gives the asserted action on $P$.
It is not necessary, and generally would not be correct, to define
the action of $\mathfrak u$ on $\mathfrak u/\mathfrak k$ by
brackets.
\end{proof}

\subsection{Trace positivity and full isotypic sign spaces}

The remaining algebraic task is to find invariant sign subspaces that
belong to a finite intrinsic list. In the infinite case we choose an
invariant submodule of the transferred dual representation. The
proposition below says why its sign parts are sums of whole isotypic
components; these sums form the required finite list.

\begin{proposition}[Pure signs in full isotypic components]
\label{alg:signs}
If $B\ne0$, let
\begin{equation}
 P'=\operatorname{soc}_{\mathfrak l}(B^*)\subset P
 \label{alg:socle}
\end{equation}
under the transferred representation of
Lemma~\ref{alg:characteristic-representation}; the socle is the sum of all
irreducible submodules. It is nonzero. Its two $H$-sign parts are
sums of full $\mathfrak u$-isotypic components, and in particular
are invariant under the full transferred $\mathfrak u$ action.

If $\mathfrak h$ is finite dimensional and $\mathfrak l=\mathfrak h$,
then $E$ is a semisimple $\mathfrak k$-module whose $H$-sign parts
are sums of full isotypic components. In either case there are only
finitely many possible sums of isotypic components in the module
under consideration.
\end{proposition}

We prove the proposition by putting the trace identity into a finite
quotient. That quotient is auxiliary: $\mathfrak l$, its representation
on $B^*$, and the socle in the statement are already intrinsically
defined. The reduction isolates the part of $H$ whose trace must vanish;
positivity will then rule out constituents containing both signs.

\subsubsection{A finite quotient and its reductive stabilizer}

We use finite-dimensional Lie theory only for the quotients and modules
specified below; the triangularization and complete-reducibility inputs
are stated in \citet[Theorems~I.3.7 and~I.5.20(b)]{Milne2013}.

\begin{lemma}
\label{alg:radical-action}
Let $g$ be a finite-dimensional complex Lie algebra with solvable
radical $r$. On every irreducible finite-dimensional $g$-module,
$r$ acts by scalars and $[g,r]$ acts trivially. On every
finite-dimensional $g$-module each element of $[g,r]$ acts
nilpotently.
\end{lemma}

\begin{proof}
Triangularize the action of $r$ by Lie's theorem. Its finite set of
diagonal characters is preserved by conjugation by
$\exp(t\rho(x))$, for $x\in g$, since $r$ is an ideal.
Dependence on $t$ is continuous, so every character is fixed;
differentiating yields $\chi([x,r])=0$ for each such character.
Choose a nonzero common eigenspace for $r$, with character $\chi$.
For a vector $v$ in this eigenspace, $y\in r$ and $x\in g$,
\[
 \rho(y)\rho(x)v
 =\rho(x)\rho(y)v+\rho([y,x])v
 =\chi(y)\rho(x)v.
\]
Thus this eigenspace is $g$-invariant. In an irreducible module it
is the whole module, proving the scalar assertion. A composition
series of a general module then makes $[g,r]$ act strictly
triangularly, proving nilpotence.
\end{proof}

By Lemma~\ref{alg:intrinsic-l}, choose an open ideal $N$ of
$\mathfrak l$ small enough that the quotient
\[
 g=\mathfrak l/N
\]
carries both $\xi|_{\mathfrak l}$ and the representation on $B$.
The latter representation is continuous: choose finitely many lifts
of a basis of $B$ and use continuity of their brackets followed by
projection to $B$. Thus the intersection of its kernel with
$\ker(\xi|_{\mathfrak l})$ contains a sufficiently small open ideal.
In the finite case one may take $N=0$. The ideal $N$ is contained
in $\mathfrak u$, since it is an ideal annihilated by $\xi$.
Every quotient in use is $H$-invariant, and its $H$-action is
semisimple with real weights.

\begin{lemma}
\label{alg:finite-trace}
For $p\in\mathfrak l_0$, identity~\eqref{alg:entropy-trace} becomes
\begin{equation}
 \begin{split}
 c\,\xi(p)
 ={}&\tr\bigl(p|(\mathfrak l/\mathfrak u)_+\bigr)
       +\tr\bigl(p|B_+\bigr)
       -\tr\bigl(p|B_-\bigr).
 \end{split}
 \label{alg:three-traces}
\end{equation}
The functional on $g$ annihilates $[g,r_g]$, where $r_g$ is the
radical of $g$.
\end{lemma}

\begin{proof}
Use the filtration $0\subset P\subset L\subset E$. Its quotients
are $P\simeq B^*$, $L/P\simeq\mathfrak l/\mathfrak u$ and
$E/L\simeq B$. A zero-weight $p$ lies in $\mathfrak k$, so acts
on all three quotients and commutes with $H$. On the positive part
of the dual module its trace is minus the trace on $B_-$.
Additivity of trace proves~\eqref{alg:three-traces}.

Lift a weight-zero element of $[g,r_g]$ to a weight-zero element
$p\in\mathfrak l$. Such a lift exists by the formal weight
projection property. It lies in $\mathfrak k\subset\mathfrak u$.
Lemma~\ref{alg:radical-action} makes its actions on $B$ and on the
adjoint representation of $g$ nilpotent. Its action on
$\mathfrak l/\mathfrak u$ is defined, precisely because
$p\in\mathfrak u$, and inherits nilpotence from the latter
adjoint action. All three traces in~\eqref{alg:three-traces}
vanish, so $\xi(p)=0$. Nonzero weights are already annihilated by
$\xi$. Weight decomposition in this finite quotient proves the
claim on all of $[g,r_g]$.
\end{proof}

Set $g'=g/[g,r_g]$. It is reductive, because its radical is central.
Denote the induced functional by $\xi'$ and its stabilizer by
$\mathfrak u'$. Since $[g,r_g]$ is an ideal annihilated by $\xi$,
its full inverse image in $\mathfrak l$ lies in $\mathfrak u$.
Consequently $\mathfrak u$ is the full inverse image of
$\mathfrak u'$ and
\begin{equation}
 \mathfrak l/\mathfrak u\simeq g'/\mathfrak u'.
 \label{alg:reductive-tangent}
\end{equation}

\begin{lemma}
\label{alg:reductive-stabilizer}
The algebra $\mathfrak u'$ is reductive. Its center on the
semisimple summand of $g'$ is toral. The restriction to
$\mathfrak u'$ of every irreducible $g'$-module is semisimple,
and $g'/\mathfrak u'$ is a semisimple $\mathfrak u'$-module.
\end{lemma}

\begin{proof}
Write $g'=z\oplus s$, with $s$ semisimple, and write $H_s$ for the
semisimple-summand component of the image of $H$. Its adjoint is
semisimple, so $H_s$ is a semisimple element of $s$. Let $A\in s$
be the Killing-dual of $\xi'|_s$. Every element of $g'$ commuting
with the image of $H$ has a zero-weight lift to
$\mathfrak l$, hence a lift in $\mathfrak k$. It therefore
stabilizes $\xi'$. Also $H$ stabilizes $\xi'$. These assertions
give
\[
 A\in Z_s(H_s),\qquad
 [A,Z_s(H_s)]=0.
\]
Thus $A$ lies in the center of the reductive centralizer of the
semisimple element $H_s$. To recall the toral-center assertion,
choose a Cartan subalgebra of $s$ containing $H_s$. Its root
decomposition describes $Z_s(H_s)$ as that Cartan subalgebra plus
the root spaces whose roots vanish on $H_s$. The resulting root
subsystem gives its semisimple derived algebra, and its center is
the subspace of the Cartan subalgebra annihilated by this subsystem.
Every element of this center is therefore semisimple. In particular
$A$ is semisimple. Applying the same centralizer description to $A$
shows that
\[
 \mathfrak u'=z\oplus Z_s(A)
\]
is reductive, with toral center on the $s$ summand.
For these root-centralizer and weight facts, see
\citet[Theorems~6.30, 6.38 and~8.2]{KirillovPreliminary}.

For an irreducible $g'$-module the global center $z$ acts by scalars.
The derived algebra of $Z_s(A)$ is semisimple, so its action is
completely reducible. The center of $Z_s(A)$ consists of commuting
semisimple elements of $s$ and hence acts simultaneously
diagonalizably in every finite-dimensional $s$-module. Splitting
into its eigenspaces and then into irreducible derived-algebra
modules proves complete reducibility over $\mathfrak u'$.
The same argument applies to the adjoint module of $g'$ and its
$\mathfrak u'$-invariant quotient $g'/\mathfrak u'$.
\end{proof}

\subsubsection{The positivity argument}

\begin{proof}[Proof of Proposition~\ref{alg:signs}]
The image of $H$ belongs to $\mathfrak u'$. Decompose it using the
center and derived algebra of this reductive algebra:
\[
 H=H_{\mathrm c}+H_{\mathrm{ss}},
 \qquad
 H_{\mathrm c}\in Z(\mathfrak u'),
 \quad H_{\mathrm{ss}}\in[\mathfrak u',\mathfrak u'].
\]
Lift $H_{\mathrm{ss}}$ to a weight-zero element
$p\in\mathfrak l$; this is possible because $H_{\mathrm{ss}}$
commutes with $H$. The functional $\xi'$ annihilates
$[\mathfrak u',\mathfrak u']$, since it is fixed by
$\mathfrak u'$, so $\xi(p)=0$.

Filter $B$ first by irreducible $g$-modules. By
Lemma~\ref{alg:radical-action} each factor is a $g'$-module, and by
Lemma~\ref{alg:reductive-stabilizer} its restriction to
$\mathfrak u'$ is semisimple. Use these irreducible
$\mathfrak u'$ constituents for the trace computation on $B$;
use the semisimple decomposition of $g'/\mathfrak u'$ for the
other quotient in~\eqref{alg:three-traces}. The filtrations are
$H$-invariant. Since $H$ is semisimple and $p$ has weight zero,
passing to either sign subspace respects all these quotient trace
computations.

On one irreducible constituent, $H_{\mathrm c}$ acts by a scalar
$a$, whereas $H_{\mathrm{ss}}$ acts by an operator $T$ of trace
zero. The scalar $a$ is real, as it equals the average of the real
eigenvalues of $H$. The operator $T=H-a\id$ is semisimple with
real eigenvalues, say $t_1,\ldots,t_d$, and
$\sum_i t_i=0$. There are no zero $H$-weights on the modules in
\eqref{alg:three-traces}. If both signs occur, the set
\[
 I_+=\{i:a+t_i>0\}
\]
is a nonempty proper upper segment of the ordered $t_i$ spectrum.
Every entry in this segment is strictly larger than every entry
outside it. Since the total average is zero, its sum is strictly
positive. If only one sign occurs, its sum is zero, with the empty
sum included. Thus
\[
 \tr(T|V_+)\ge0,
 \qquad
 \tr(T|V_+)-\tr(T|V_-)=2\tr(T|V_+)\ge0,
\]
with strict inequality in each expression when $V$ has both signs.

Evaluate~\eqref{alg:three-traces} at the chosen lift of
$H_{\mathrm{ss}}$. Its left side is zero, and all the constituent
contributions just described are nonnegative. None can be strictly
positive. Every irreducible constituent in these computations is
therefore pure sign. Moreover isomorphic $\mathfrak u'$ modules
have the same spectrum for the fixed element $H$, so they have the
same sign.

The finite-dimensional nonzero module $B^*$ has an irreducible
submodule, so its socle $P'$ is nonzero. It is a semisimple
$g$-module. Its irreducible constituents are duals of irreducible
quotients of $B$, hence occur among the duals of the factors just
considered. Their restrictions to $\mathfrak u'$ are semisimple,
with pure-sign irreducible constituents by the same result and with
signs reversed under duality.
Isomorphic constituents cannot be assigned opposite signs.
Consequently each sign part of $P'$ is a sum of full
$\mathfrak u$-isotypic components, for the transferred dual
representation. This proves the assertion of full
$\mathfrak u$ invariance.

In the finite case $B=0$, $\mathfrak u=\mathfrak k$, and
\eqref{alg:reductive-tangent} identifies $E$ with
$g'/\mathfrak u'$. The same constituent argument gives its
semisimple isotropy representation and pure-sign isotypic
decomposition. Any finite-dimensional semisimple module has only
finitely many distinct isotypic components and therefore finitely
many sums of entire components.
\end{proof}

\subsection{Intrinsic fields, real structures, and axis coherence}

The finite quotient used for the proof is auxiliary. The subalgebra
$\mathfrak l$ is intrinsic by~\eqref{alg:equicontinuity};
$L,P,B,\mathfrak u$ and the representation transferred to $P$ are
intrinsic from their definitions. The socle in~\eqref{alg:socle}
is also intrinsic. In particular none of these objects depends on a
choice of coordinates in the minimal ideals or on the finite
quotient $g$.

\begin{proposition}
\label{alg:axis-coherence}
The planes and subalgebras above are measurable on a common conull
set and commute with the dynamics and real conjugation. Their
invariant homogeneous formal fields have full ambient jets smooth
and holonomic along both entire typical strong axes, with the weak
versions obtained by flow transport.

In the infinite case $P'$ is flip-invariant. Its actual stable and
unstable parts
\[
 P^s=P'\cap\Es,
 \qquad P^u=P'\cap\Eu
\]
are both nonzero, have the same rank almost everywhere, and have
the stated coherent formal extensions. They remain invariant under
the full transferred $\mathfrak u$ action. In the finite case the
formal sign fields have center values exactly $\Es$ and $\Eu$.
All these formal fields are invariant under the signed dilation.
\end{proposition}

\begin{proof}
We first explain measurability without choosing a measurable
classification of minimal ideals. Replace the countably many
completed-measurable background coefficients by Borel versions on a
common conull Borel set, refined under the countable group generated
by $f$ and flip. For this purpose, coefficient
sequences are coded in the countable product of the usual standard
Borel space $\C$. This measurable structure is distinct from the
formal linear topology with discrete scalars used in the algebraic
arguments. The symmetry algebra is given by
countably many linear jet equations. For the intrinsic condition
\eqref{alg:equicontinuity}, fix an output jet order. It asks for an
input jet order such that every power of the adjoint annihilates the
corresponding input kernel after output projection. For each power
this is a finite-jet linear inclusion test; the full formal solution
images are obtainable by the stabilized finite linear tests from
Lemma~\ref{sym:full-projections}. The quantifiers over orders and powers
are countable. Thus membership in $\mathfrak l$ is a Borel
relation between the base point and coefficient sequences. Its
finite-dimensional projections are analytic, hence universally
measurable \citep[Proposition~12.2 and Corollary~14.10]{Tserunyan2025}.
To recover the tangent plane measurably, for each $d$ code $d$
independent projected vectors. Their analytic relation has domain
$A_d$, the set where the dimension is at least $d$, and admits a
universally measurable tuple selection
\citep[Section~4E.9]{Moschovakis2009}. Restrict this already chosen
tuple to $A_d\setminus A_{d+1}$ and take its span. This gives the
Grassmannian-valued plane for the completed Liouville measure;
the exact-dimension strata need not be analytic. Subsequent countable
coefficient choices can again be replaced by Borel versions before
further coding. No Borel selector, or common Borel refinement invariant
under every real flow time, is required.

Once the finite-dimensional plane $L$ is known, brackets and finite
jet lifts determine the representation on $B$. The socle of its
dual is specified by finite-dimensional invariant-subspace and rank
conditions. Indeed invariance is a closed Grassmannian incidence
condition for a finite basis of the represented matrix algebra;
reducibility is a finite union of projections along compact nested
Grassmannians. Thus irreducibility is Borel, and the preceding
independent-tuple procedure selects the span of the irreducible
submodules. These operations give measurable planes and, where
needed, measurable choices of finite jets. The homogeneous formal
extensions are computed to each requested order from finitely many
such jet operations. All resulting constructions are natural for
the dynamics. Their dimensions are therefore almost everywhere
constant by ergodicity.

Conjugation preserves the intrinsic definitions and the real
$H$-sign spaces; hence the complex planes descend to real ones.
Flip intertwines the two affine-form data, preserves the
underlying symmetry structure up to the harmless sign of the
contact form, and consequently preserves $P'$. It exchanges the
actual stable and unstable parts. Since it preserves smooth volume
and these two ranks are almost everywhere constant, the ranks are
equal. Their sum is the positive rank of $P'$, so both are nonzero.

For smoothness and coherence, use the formal identifications along
whole axes supplied by Proposition~\ref{sym:axis-transport}.
They carry $\mathfrak h,\mathfrak k,\xi$, hence all the intrinsic
objects just constructed, smoothly along the axis. The transported
plane $P'$ is preserved by the projections to the two actual sign
bundles at typical points. These projections are continuous and
uniformly transverse. By density they preserve the smooth extension
of $P'$ everywhere on the axis patch. Their restrictions are
continuous idempotents, so their ranks are locally constant; their
images give continuous sign subbundles.

In this trivialization, Proposition~\ref{alg:signs} places each
sign subspace, at typical points, in a fixed finite list of full
isotypic sums. The finite list is closed, so the same holds at all
points by density. A continuous choice in that list is locally
constant. The same reasoning applies to the finite-case sign
spaces in $E$. Proposition~\ref{sym:axis-transport}, including
its version for local smooth families of paths, now gives smooth
full formal jets and their holonomicity along the whole axes.
The homogeneous extensions are invariant under $\mathfrak h$ by
Proposition~\ref{sym:homogeneous-fields}, and thus under $H$.
Their center values are the prescribed actual sign planes. Finally,
flow shifts intertwine all these intrinsic objects; the section
fields and their horizontal lifts can be taken time-independent
within a flow box. This supplies the weak-axis extensions with the
required agreement at regular translates.
\end{proof}

\begin{theorem}[Algebraic dichotomy]
\label{alg:dichotomy}
At typical centers exactly one of the following alternatives holds,
with the same alternative almost everywhere.
\begin{enumerate}[label=(\roman*)]
\item The symmetry algebra is finite dimensional. There are
$\mathfrak h$-invariant formal fields with center values $\Es$
and $\Eu$, obtained from finite isotypic sign choices. Their full
ambient jets are smooth and holonomic along both entire typical
axes and equivariant under the dynamics.
\item The symmetry algebra is infinite dimensional. The intrinsic
proper open subalgebra $\mathfrak l\supset\mathfrak k$ determines
an involutive coisotropic formal distribution $L$, whose
characteristic distribution is $P=L^\omega$. It contains an
intrinsic nonzero plane $P'$ with nonzero actual stable and unstable
parts $P^s,P^u$. These parts are invariant under the full
$\mathfrak u$ action transferred from $B^*$, and all these formal
fields have the same whole-axis coherence and equivariance.
\end{enumerate}
\end{theorem}

\begin{proof}
The algebraic assertions are
Propositions~\ref{alg:coisotropic} and~\ref{alg:signs}; the real,
measurable and coherent conclusions are
Proposition~\ref{alg:axis-coherence}. The subalgebra
$\mathfrak l$ gives an involutive homogeneous formal distribution
by Proposition~\ref{sym:homogeneous-fields}. Coisotropy gives its
characteristic distribution. Invariance and ergodicity make the
alternative almost everywhere constant.
\end{proof}

We record the precise representation-theoretic input to the
characteristic-leaf argument in the next section. First integrals
of $L$ define a formal base with tangent space $B$. For
$p\in\mathfrak l$, the projected symmetry fixes the base origin;
the negative of its linearization is the bracket action on $B$.
The cotangent action, under Hamiltonian translation frames for the
characteristic leaf, is exactly the transferred dual action on $P$.
Elements of $\mathfrak u$ have values spanning $P$ and preserve
that leaf. Proposition~\ref{alg:signs} therefore supplies invariant
initial subspaces for this full cotangent action. This is the
invariance needed to prove that $P^s$ and $P^u$ are parallel for
the canonical characteristic affine connection. The construction
of that connection, its restriction to actual plaques, and the
realization of mixed rectangles are carried out below; no
convergence of formal flows at nonzero parameters is asserted here.
